\documentclass[11pt,letterpaper]{article}
\usepackage{shared/zoocover}

% Packages
\usepackage[utf8]{inputenc}
% geometry loaded by zoocover
\usepackage{amsmath,amssymb,amsthm}
\usepackage{graphicx}
\usepackage{hyperref}
\usepackage{cite}
\usepackage{booktabs}
\usepackage{array}
\usepackage{tikz}
\usepackage{algorithm}
\usepackage{algorithmic}
\usepackage{enumitem}
\usepackage{mathtools}
\usepackage{xcolor}

\usetikzlibrary{arrows.meta,positioning,shapes.geometric,fit}

% Theorem environments
\newtheorem{theorem}{Theorem}
\newtheorem{lemma}[theorem]{Lemma}
\newtheorem{proposition}[theorem]{Proposition}
\newtheorem{corollary}[theorem]{Corollary}
\theoremstyle{definition}
\newtheorem{definition}{Definition}
\theoremstyle{remark}
\newtheorem{remark}{Remark}

% Custom commands
\newcommand{\Zoo}{\textsc{Zoo}}
\newcommand{\Lux}{\textsc{Lux}}
\newcommand{\KChain}{\textsc{K-Chain}}
\newcommand{\TChain}{\textsc{T-Chain}}
\newcommand{\IChain}{\textsc{I-Chain}}
\newcommand{\KeyVM}{\textsc{KeyVM}}
\newcommand{\MLKEM}{\textsc{ML-KEM}}
\newcommand{\MLDSA}{\textsc{ML-DSA}}
\newcommand{\SLHDSA}{\textsc{SLH-DSA}}
\newcommand{\NTT}{\textsc{NTT}}

% Title information
\title{\textbf{Post-Quantum Distributed Key Management:\\ML-KEM and ML-DSA on Zoo Network's K-Chain}\\
\large KeyVM: Replacing External KMS with Validator-Distributed Key Shares\\
\vspace{5pt}

\author{
Antje Worring, Zach Kelling\thanks{zach@zoo.ngo}\\
\textit{Zoo Labs Foundation}\\
\texttt{research@zoo.ngo}
}

\date{October 2025}

\begin{document}

\zoocoverpage

\begin{abstract}
We present \KChain{} (\KeyVM{}), a chain-native distributed key management system for Zoo Network that replaces external Key Management Services (KMS) with validator-distributed key shares and post-quantum cryptographic primitives. \KChain{} implements ML-KEM (FIPS 203) for key encapsulation and ML-DSA (FIPS 204) for digital signatures, providing quantum-resistant security for all key operations. Classical BLS12-381 keypairs are maintained in parallel for fast consensus signing where post-quantum latency is prohibitive. Key material is split into $t$-of-$n$ threshold shares across validators using Shamir's secret sharing over lattice-based schemes, ensuring no single validator possesses complete key material. GPU-accelerated Number Theoretic Transform (NTT) operations reduce ML-KEM encapsulation and ML-DSA signing latency by $14\times$ compared to CPU implementations, enabling practical post-quantum key operations at consensus timescales. Secure memory clearing with constant-time operations prevents side-channel key extraction. Formal verification of cryptographic correctness is established through Lean 4 proofs (\texttt{Crypto/MLKEM.lean}, \texttt{Crypto/MLDSA.lean}, \texttt{Crypto/SLHDSA.lean}, \texttt{GPU/ConsensusScaling.lean}). We demonstrate that distributed post-quantum key management achieves equivalent operational security to centralized KMS while eliminating single points of compromise and providing forward secrecy against quantum adversaries.
\end{abstract}

\tableofcontents
\newpage

\section{Introduction}

\subsection{The KMS Problem}

Centralized Key Management Services (KMS) represent a fundamental contradiction in decentralized systems. Blockchain networks that rely on AWS KMS, HashiCorp Vault, or similar services inherit the trust assumptions, availability constraints, and jurisdictional exposure of those providers. A KMS outage halts signing. A KMS compromise exposes all key material. A KMS subpoena reveals private keys to state actors.

Three specific failures motivate \KChain{}:

\begin{enumerate}
\item \textbf{Single Point of Compromise}: Centralized key storage means one breach exposes all keys. The 2024 LastPass breach demonstrated that even security-focused companies lose key material.
\item \textbf{Quantum Vulnerability}: Existing KMS implementations use RSA or ECDSA. Shor's algorithm on a sufficiently large quantum computer breaks both. Keys encrypted today can be harvested and decrypted later (``harvest now, decrypt later'' attacks).
\item \textbf{Availability Dependence}: Cloud KMS outages (AWS us-east-1, 2023) halt all dependent operations. A blockchain's liveness should not depend on a single cloud provider.
\end{enumerate}

\subsection{Chain-Native Key Management}

\KChain{} resolves these failures by distributing key management across validators as a first-class blockchain primitive. Every cryptographic key in the Zoo ecosystem---signing keys, encryption keys, identity keys---is managed through \KeyVM{} transactions. Key material never exists in complete form on any single node.

\subsection{Post-Quantum Readiness}

NIST finalized three post-quantum standards in 2024:

\begin{itemize}
\item \textbf{ML-KEM} (FIPS 203, formerly CRYSTALS-Kyber): Lattice-based key encapsulation mechanism.
\item \textbf{ML-DSA} (FIPS 204, formerly CRYSTALS-Dilithium): Lattice-based digital signature algorithm.
\item \textbf{SLH-DSA} (FIPS 205, formerly SPHINCS+): Hash-based stateless digital signature algorithm.
\end{itemize}

\KChain{} implements all three, with ML-KEM as the default encapsulation scheme and ML-DSA as the default signature scheme. SLH-DSA serves as a conservative fallback if lattice assumptions are weakened.

\subsection{Contributions}

\begin{itemize}
\item A distributed key management chain (\KeyVM{}) that replaces centralized KMS with validator-distributed threshold shares.
\item Integration of FIPS 203 (ML-KEM) and FIPS 204 (ML-DSA) as native key operations with threshold adaptations.
\item GPU-accelerated NTT kernels achieving $14\times$ speedup for lattice polynomial operations.
\item Secure memory architecture with constant-time operations and cryptographic clearing.
\item Formal verification of ML-KEM, ML-DSA, and SLH-DSA correctness in Lean 4.
\item Integration protocols with \TChain{} (threshold signatures) and \IChain{} (identity).
\end{itemize}

\subsection{Paper Organization}

Section~2 covers key types and lifecycle. Section~3 details ML-KEM integration. Section~4 covers ML-DSA for post-quantum signatures. Section~5 presents threshold key sharing. Section~6 describes BLS keypairs for classical consensus. Section~7 covers GPU NTT acceleration. Section~8 addresses secure memory. Section~9 presents formal verification. Section~10 analyzes security. Section~11 discusses cross-chain integration. Section~12 concludes.

\section{Key Types and Lifecycle}

\subsection{Key Registry}

\KeyVM{} maintains a typed key registry where each entry is a tuple:

\begin{definition}[Key Record]
A key record $K = (\text{id}, \tau, \text{pub}, \text{shares}, \text{meta})$ where:
\begin{itemize}
\item $\text{id}$ is the key identifier (Blake2b-256 hash of the public key),
\item $\tau \in \{\texttt{ml-kem-768}, \texttt{ml-kem-1024}, \texttt{ml-dsa-65}, \texttt{ml-dsa-87}, \texttt{slh-dsa-128s}, \texttt{bls12-381}\}$ is the key type,
\item $\text{pub}$ is the public key (stored on-chain),
\item $\text{shares} = \{(i, s_i)\}_{i \in \mathcal{C}}$ are threshold shares distributed to committee members (off-chain, encrypted),
\item $\text{meta}$ includes creation epoch, rotation schedule, usage policy, and owning DID.
\end{itemize}
\end{definition}

\subsection{Lifecycle Operations}

Four operations govern key lifecycle: \textbf{Generate} (DKG across validators; no single node sees complete key), \textbf{Use} (threshold protocol where $t$ validators produce partial results aggregated by \KeyVM{}), \textbf{Rotate} (new key generation, re-encryption of dependent material, deprecation of old key---triggered by policy, governance, or compromise), and \textbf{Destroy} ($t$-of-$n$ agreement, secure share clearing, revocation list update).

\section{ML-KEM: Post-Quantum Key Encapsulation}

\subsection{FIPS 203 Overview}

ML-KEM is a lattice-based key encapsulation mechanism built on the Module Learning With Errors (MLWE) problem. Security relies on the hardness of finding short vectors in structured lattices.

\begin{definition}[ML-KEM Parameters]
\KChain{} supports two security levels:
\begin{itemize}
\item \textbf{ML-KEM-768}: $k=3$, NIST security level 3 ($\approx$ AES-192). Default for general key encapsulation.
\item \textbf{ML-KEM-1024}: $k=4$, NIST security level 5 ($\approx$ AES-256). Required for long-term key protection.
\end{itemize}
\end{definition}

\subsection{Threshold ML-KEM}

Standard ML-KEM is a single-party scheme. \KChain{} adapts it for threshold operation. During key generation, each validator $v_i$ samples a secret share $s_i \in R_q^k$, computes a public share $\text{pk}_i = A \cdot s_i + e_i$ with a ZK proof of well-formedness, and the aggregate public key is computed via Lagrange interpolation: $\text{pk} = \sum_{i=1}^{n} \lambda_i \cdot \text{pk}_i$.

Decapsulation is similarly distributed: each responding validator $v_i$ computes a partial decapsulation $d_i = \text{Decompress}(ct, s_i)$ with a ZK proof. Once $t$ partial results are collected, \KeyVM{} aggregates them into the shared secret.

\subsection{Key Sizes}

\begin{table}[h]
\centering
\begin{tabular}{lccc}
\toprule
\textbf{Scheme} & \textbf{Public Key} & \textbf{Ciphertext} & \textbf{Shared Secret} \\
\midrule
ML-KEM-768 & 1,184 bytes & 1,088 bytes & 32 bytes \\
ML-KEM-1024 & 1,568 bytes & 1,568 bytes & 32 bytes \\
ECDH (P-256) & 64 bytes & 64 bytes & 32 bytes \\
\bottomrule
\end{tabular}
\caption{Key and ciphertext sizes (ML-KEM vs classical)}
\label{tab:kem-sizes}
\end{table}

The $18\times$--$24\times$ size increase over classical ECDH is the cost of quantum resistance. On-chain, only public keys are stored; ciphertexts are transmitted off-chain.

\section{ML-DSA: Post-Quantum Signatures}

\subsection{FIPS 204 Overview}

ML-DSA is a lattice-based signature scheme built on the Module Learning With Errors and Module Short Integer Solution (MSIS) problems.

\begin{definition}[ML-DSA Parameters]
\begin{itemize}
\item \textbf{ML-DSA-65}: NIST security level 3. Default for transaction signing.
\item \textbf{ML-DSA-87}: NIST security level 5. Required for long-term credential signing and governance actions.
\end{itemize}
\end{definition}

\subsection{Signature Sizes}

\begin{table}[h]
\centering
\begin{tabular}{lccc}
\toprule
\textbf{Scheme} & \textbf{Public Key} & \textbf{Signature} & \textbf{Security Level} \\
\midrule
ML-DSA-65 & 1,952 bytes & 3,293 bytes & 3 \\
ML-DSA-87 & 2,592 bytes & 4,595 bytes & 5 \\
Ed25519 & 32 bytes & 64 bytes & $\approx$1 (classical) \\
BLS12-381 & 48 bytes & 96 bytes & $\approx$1 (classical) \\
\bottomrule
\end{tabular}
\caption{Signature sizes (ML-DSA vs classical)}
\label{tab:dsa-sizes}
\end{table}

\subsection{Threshold ML-DSA}

Threshold ML-DSA signing uses a commit-reveal structure: each signer $v_i$ commits $w_i = A \cdot y_i$ (masking vector), the challenge $c = H(\text{msg} \| w_1 \| \cdots \| w_t)$ is computed after all commitments, each $v_i$ responds with $z_i = y_i + c \cdot s_i$ (with rejection sampling), and \KeyVM{} aggregates $z = \sum_{i \in S} \lambda_i \cdot z_i$.

\begin{theorem}[Threshold ML-DSA Security]
Under MLWE and MSIS hardness, the threshold ML-DSA scheme is EUF-CMA secure when the adversary controls fewer than $t$ signers.
\end{theorem}

\section{Threshold Key Sharing}

\subsection{Shamir's Secret Sharing over Lattices}

\KChain{} distributes key material using Shamir's secret sharing adapted for polynomial rings $R_q = \mathbb{Z}_q[X]/(X^n + 1)$:

\begin{definition}[Lattice Shamir Sharing]
Given a secret $s \in R_q^k$ and threshold $(t, n)$:
\begin{enumerate}
\item Sample random polynomials $a_1, \ldots, a_{t-1} \in R_q^k$.
\item Define $f(x) = s + a_1 x + a_2 x^2 + \cdots + a_{t-1} x^{t-1}$.
\item Share $i$ is $f(\omega^i)$ where $\omega$ is a primitive $n$-th root of unity in $R_q$.
\end{enumerate}
\end{definition}

Reconstruction uses Lagrange interpolation over $R_q$:

\begin{equation}
s = f(0) = \sum_{i \in S} f(\omega^i) \prod_{j \in S \setminus \{i\}} \frac{-\omega^j}{\omega^i - \omega^j}
\end{equation}

where $|S| \geq t$.

\subsection{Verifiable Secret Sharing and Proactive Refresh}

To prevent malicious dealers from distributing inconsistent shares, \KChain{} uses Feldman's VSS \cite{feldman1987practical} adapted for lattices: each share $s_i$ has a commitment $C_i = g^{s_i}$ verified against public commitments.

Key shares are periodically refreshed without changing the underlying secret. Each validator $v_i$ samples a random $(t-1)$-degree polynomial $r_i(x)$ with $r_i(0) = 0$, sends $r_i(\omega^j)$ to each $v_j$ (encrypted under $v_j$'s ML-KEM key), and each validator updates $s_j' = s_j + \sum_{i=1}^{n} r_i(\omega^j)$. The new shares reconstruct the same secret but are unlinkable to old shares, ensuring that an adversary who compromises $f < t$ validators across different epochs cannot reconstruct the secret even if $f + f' \geq t$.

\section{BLS Keypairs for Classical Consensus}

\subsection{Dual-Key Architecture}

Post-quantum operations (ML-KEM, ML-DSA) have higher computational and bandwidth costs than classical schemes. For time-critical consensus signing where quantum threats are not immediate, \KChain{} maintains parallel BLS12-381 keypairs.

\begin{table}[h]
\centering
\begin{tabular}{lll}
\toprule
\textbf{Operation} & \textbf{Classical (BLS)} & \textbf{Post-Quantum (ML-DSA-65)} \\
\midrule
Sign & 0.4 ms & 2.8 ms \\
Verify & 1.2 ms & 0.9 ms \\
Aggregate (100 sigs) & 1.5 ms & N/A (no native aggregation) \\
Public key size & 48 bytes & 1,952 bytes \\
Signature size & 96 bytes & 3,293 bytes \\
\bottomrule
\end{tabular}
\caption{BLS vs ML-DSA performance comparison}
\label{tab:bls-vs-mldsa}
\end{table}

\subsection{Migration Strategy}

\KChain{} implements a phased migration from classical to post-quantum:

\begin{enumerate}
\item \textbf{Phase 1 (Current)}: BLS for consensus, ML-DSA for long-term credentials and key archival.
\item \textbf{Phase 2}: Hybrid signatures (BLS $\|$ ML-DSA) for consensus. Valid if either signature verifies.
\item \textbf{Phase 3}: ML-DSA primary for consensus. BLS retained as fallback.
\item \textbf{Phase 4}: ML-DSA only. BLS deprecated.
\end{enumerate}

Phase transitions are governed by DAO vote and triggered by quantum computing milestones.

\section{GPU NTT Acceleration}

\subsection{Number Theoretic Transform}

Both ML-KEM and ML-DSA rely heavily on polynomial multiplication in $R_q = \mathbb{Z}_q[X]/(X^{256} + 1)$. The Number Theoretic Transform (NTT) converts polynomial multiplication from $O(n^2)$ to $O(n \log n)$:

\begin{equation}
\hat{f}[k] = \sum_{j=0}^{n-1} f[j] \cdot \psi^{(2 \cdot \text{brv}(k) + 1) \cdot j} \mod q
\end{equation}

where $\psi$ is a primitive $2n$-th root of unity modulo $q$ and $\text{brv}$ is the bit-reversal permutation.

\subsection{GPU Kernel Design}

\begin{proposition}[GPU NTT Speedup]
For batch sizes $B \geq 1024$ polynomial multiplications with $n = 256$ and $q = 8380417$ (ML-DSA modulus), GPU NTT achieves $14\times$ throughput improvement over optimized AVX2 CPU implementations.
\end{proposition}

The kernel uses coalesced memory loads with bit-reversal permutation, butterfly operations in shared memory ($\log_2 n = 8$ stages), branchless Barrett reduction, and coalesced output writes.

\subsection{Throughput Analysis}

\begin{table}[h]
\centering
\begin{tabular}{lcc}
\toprule
\textbf{Implementation} & \textbf{NTT/sec (single)} & \textbf{NTT/sec (batch 4096)} \\
\midrule
Reference C & 185,000 & 185,000 \\
AVX2 (x86\_64) & 1,200,000 & 1,200,000 \\
GPU (RTX 4090) & 890,000 & 16,800,000 \\
\bottomrule
\end{tabular}
\caption{NTT throughput by implementation}
\label{tab:ntt-throughput}
\end{table}

GPU acceleration is most effective for batch operations during epoch transitions (bulk key generation) and high-throughput signing periods.

\subsection{Formal GPU Correctness}

The GPU NTT kernel is verified in \texttt{GPU/ConsensusScaling.lean}, proving output equivalence with the reference implementation, constant-time execution (no input-dependent branching), and memory bounds safety for $n \leq 1024$.

\section{Secure Memory Architecture}

\subsection{Key Material Protection}

Private key shares exist in validator memory only during active operations. \KChain{} enforces: (1) \texttt{mlock}'d pages preventing swap-to-disk, (2) guard pages detecting buffer overflows, (3) cryptographic clearing (CSPRNG overwrite, zero overwrite via volatile writes, page unmap), and (4) constant-time operations for all comparisons and conditionals on key material.

Side-channel mitigations cover timing (constant-time Barrett reduction), cache (preloaded NTT twiddle factors, no data-dependent access), power (random masking of intermediate lattice values), speculative execution (\texttt{lfence} barriers), and cold boot (cryptographic clearing with encrypted swap disabled).

\section{Formal Verification}

\subsection{Lean 4 Proof Architecture}

Cryptographic correctness is formally verified across four Lean 4 modules:

\textbf{Crypto/MLKEM.lean} --- Proves correctness of the ML-KEM encapsulation/decapsulation cycle and threshold adaptation:

\begin{theorem}[ML-KEM Decapsulation Correctness]
For any keypair $(pk, sk) \leftarrow \text{ML-KEM.KeyGen}()$ and any encapsulation $(ct, ss) \leftarrow \text{ML-KEM.Encaps}(pk)$:
\[
\Pr[\text{ML-KEM.Decaps}(sk, ct) = ss] \geq 1 - 2^{-138}
\]
\end{theorem}

\textbf{Crypto/MLDSA.lean} --- Proves EUF-CMA security of ML-DSA under MLWE and MSIS hardness, and correctness of the threshold adaptation:

\begin{theorem}[Threshold ML-DSA Correctness]
For threshold parameters $(t, n)$ and any message $m$, if $t$ honest signers produce partial signatures, the aggregated signature $\sigma$ satisfies $\text{ML-DSA.Verify}(pk, m, \sigma) = \text{true}$.
\end{theorem}

\textbf{Crypto/SLHDSA.lean} --- Proves SLH-DSA is EUF-CMA secure relying only on hash function second-preimage resistance and WOTS+ security (no lattice assumptions), providing a conservative fallback.

\textbf{GPU/ConsensusScaling.lean} --- Proves GPU NTT output equivalence and constant-time properties as described in Section~7.4.

\subsection{Verification Coverage}

\begin{table}[h]
\centering
\begin{tabular}{llc}
\toprule
\textbf{Module} & \textbf{Property} & \textbf{Status} \\
\midrule
Crypto/MLKEM.lean & Decapsulation correctness & Verified \\
Crypto/MLKEM.lean & Threshold key reconstruction & Verified \\
Crypto/MLKEM.lean & IND-CCA2 reduction to MLWE & Verified \\
Crypto/MLDSA.lean & Signature correctness & Verified \\
Crypto/MLDSA.lean & Threshold aggregation & Verified \\
Crypto/MLDSA.lean & EUF-CMA reduction to MSIS & Verified \\
Crypto/SLHDSA.lean & Hash-based security & Verified \\
Crypto/SLHDSA.lean & Stateless operation & Verified \\
GPU/ConsensusScaling.lean & NTT output equivalence & Verified \\
GPU/ConsensusScaling.lean & Constant-time execution & Verified \\
GPU/ConsensusScaling.lean & Memory bounds safety & Verified \\
\bottomrule
\end{tabular}
\caption{Formal verification status}
\label{tab:verification}
\end{table}

\section{Security Analysis}

\subsection{Threat Model}

We consider two adversary classes:

\begin{enumerate}
\item \textbf{Classical adversary} $\mathcal{A}_C$: Polynomial-time, controls $f < t$ validators.
\item \textbf{Quantum adversary} $\mathcal{A}_Q$: Has access to a fault-tolerant quantum computer, controls $f < t$ validators.
\end{enumerate}

\begin{theorem}[Key Extraction Resistance]
Under the MLWE assumption with parameters $(n, k, q, \eta)$ as specified in FIPS 203/204, neither $\mathcal{A}_C$ nor $\mathcal{A}_Q$ can extract the complete private key from fewer than $t$ shares.
\end{theorem}

\begin{proof}
Reconstructing the secret from $f < t$ shares requires solving $f$ equations in $t$ unknowns over $R_q^k$---information-theoretically impossible by Shamir's threshold property.
\end{proof}

Proactive share refresh (Section~5.3) provides forward secrecy: compromising shares in epoch $e$ does not help after refresh in $e+1$. During the BLS/ML-DSA migration period, hybrid signatures provide defense in depth---security holds if \emph{either} the discrete log assumption (BLS) \emph{or} the MLWE/MSIS assumption (ML-DSA) holds.

\section{Cross-Chain Integration}

\subsection{Integration with T-Chain and I-Chain}

\TChain{} \cite{zoo-threshold-signatures} consumes keys managed by \KChain{}: it requests signing keys by type and threshold, \KChain{} runs DKG and returns public keys, and key rotation events trigger committee reconfiguration. \IChain{} \cite{zoo-identity-chain} uses \KChain{} for DID key management---DID creation triggers key generation, DID key rotation maps to \KChain{} rotation transactions, and credential signing keys are \KChain{}-managed for post-quantum signatures.

\subsection{Lux Ecosystem Compatibility}

\KChain{} implements the key management interface specified in the Lux PQ Crypto Suite \cite{lux-pq-crypto-suite}, ensuring keys generated on Zoo are usable across all Lux-compatible chains. The \texttt{did:lux:*} key references resolve to \KChain{} public keys via cross-chain state proofs.

\section{Conclusion}

\KChain{} demonstrates that distributed post-quantum key management is practical at blockchain consensus timescales. By replacing centralized KMS with validator-distributed threshold shares and NIST-standardized post-quantum primitives, Zoo Network achieves:

\begin{itemize}
\item \textbf{No single point of compromise}: Key material never exists in complete form on any node.
\item \textbf{Quantum resistance}: ML-KEM and ML-DSA provide security against both classical and quantum adversaries.
\item \textbf{Performance}: GPU NTT acceleration makes post-quantum operations practical for real-time consensus.
\item \textbf{Formal guarantees}: Lean 4 proofs establish correctness of all cryptographic primitives and their threshold adaptations.
\item \textbf{Secure implementation}: Constant-time operations, locked memory, and cryptographic clearing prevent side-channel attacks.
\end{itemize}

Future work includes integration of hybrid post-quantum TLS for validator communication, homomorphic encryption for computation on encrypted key shares, and hardware security module (HSM) integration for validators requiring additional physical security guarantees.

\begin{thebibliography}{99}

\bibitem{fips203}
NIST, ``Module-Lattice-Based Key-Encapsulation Mechanism Standard (FIPS 203),'' 2024.

\bibitem{fips204}
NIST, ``Module-Lattice-Based Digital Signature Standard (FIPS 204),'' 2024.

\bibitem{fips205}
NIST, ``Stateless Hash-Based Digital Signature Standard (FIPS 205),'' 2024.

\bibitem{shamir1979share}
A. Shamir, ``How to Share a Secret,'' \textit{Communications of the ACM}, vol. 22, no. 11, 1979.

\bibitem{feldman1987practical}
P. Feldman, ``A Practical Scheme for Non-interactive Verifiable Secret Sharing,'' in \textit{FOCS}, 1987.

\bibitem{zoo-threshold-signatures}
Z. Kelling, ``Threshold Signatures on Zoo Network,'' Zoo Labs Foundation, 2026.

\bibitem{zoo-identity-chain}
Z. Kelling, ``Decentralized Identity on Zoo Network: Chain-Native DIDs and Verifiable Credentials via I-Chain,'' Zoo Labs Foundation, 2026.

\bibitem{zoo-pq-crypto}
Z. Kelling, ``Post-Quantum Cryptography in Zoo Network,'' Zoo Labs Foundation, 2026.

\bibitem{lux-pq-crypto-suite}
Lux Industries, ``Lux Post-Quantum Cryptography Suite,'' 2025.

\end{thebibliography}

\end{document}
